\SourceLeaf{21}
\begin{center}\Large SUMMATION\end{center}

\begin{center}\large EINIGER BESONDERER REIHEN.\end{center}

\BookRule{22mm}

\Article{1}

Unter den bedeutenderen Wahrheiten, zu denen die Theorie der Kreisteilung den Zugang eröffnet hat, nimmt die in den Disquiss. Arithmet., Art. 356, angegebene Summation einen hervorragenden Platz ein. Das liegt sowohl an ihrer besonderen Eleganz und erstaunlichen Fruchtbarkeit, die eine andere Untersuchung später ausführlicher darzustellen Gelegenheit geben wird, als auch an den ungewöhnlichen Schwierigkeiten ihres strengen und vollständigen Beweises.\quad Diese Schwierigkeiten hätte man umso weniger erwarten sollen, als sie weniger den Satz selbst als vielmehr eine gewisse Einschränkung des Satzes betreffen: Lässt man diese außer Betracht, so steht sofort ein Beweis zur Verfügung, der sich sehr leicht aus der in jenem Werk entwickelten Theorie ergibt.\quad Der Satz wurde dort in folgender Form angegeben.\quad Es sei $n$ eine Primzahl; alle quadratischen Reste von $n$ zwischen den Grenzen $1$ und $n-1$ einschließlich seien unbestimmt durch $a$ bezeichnet und alle Nichtreste zwischen denselben Grenzen durch $b$, schließlich bezeichne $\omega$ den Bogen $\frac{360^{0}}{n}$ und $k$ eine beliebige feste ganze Zahl, die nicht durch $n$ teilbar ist. Dann gilt

I.\quad für einen Wert von $n$ von der Form $4m+1$:

\[
\begin{aligned}
\Sigma\cos ak\omega&=-\tfrac12\pm\tfrac12\surd n\\
\Sigma\cos bk\omega&=-\tfrac12\mp\tfrac12\surd n,\ \text{und somit}\\
\Sigma\cos ak\omega-\Sigma\cos bk\omega&=\pm\surd n\\
\Sigma\sin ak\omega&=0\\
\Sigma\sin bk\omega&=0
\end{aligned}
\]

\SourceLeaf{22}
II.\quad für einen Wert von $n$ von der Form $4m+3$:

\[
\begin{aligned}
\Sigma\cos ak\omega&=-\tfrac12\\
\Sigma\cos bk\omega&=-\tfrac12\\
\Sigma\sin ak\omega&=\pm\tfrac12\surd n\\
\Sigma\sin bk\omega&=\mp\tfrac12\surd n\\
\Sigma\sin ak\omega-\Sigma\sin bk\omega&=\pm\surd n
\end{aligned}
\]

Diese Summationen sind an der angeführten Stelle mit voller Strenge bewiesen worden. Hier bleibt nur die Schwierigkeit, das \emph{Vorzeichen} zu bestimmen, das vor den Wurzelausdruck zu setzen ist.\quad Ohne Mühe lässt sich zwar zeigen, dass dieses Vorzeichen insofern von der Zahl $k$ abhängt, als für alle Werte von $k$, die quadratische Reste von $n$ sind, stets \emph{dasselbe} Vorzeichen gelten muss und für alle Werte von $k$, die quadratische Nichtreste von $n$ sind, das entgegengesetzte.\quad Damit kommt es nur auf den Wert $k=1$ an: Sobald das für diesen Wert geltende Vorzeichen bekannt ist, stehen auch die Vorzeichen für alle übrigen Werte von $k$ sofort zur Verfügung.\quad Doch gerade bei dieser Frage, die auf den ersten Blick zu den leichteren zu gehören scheint, stoßen wir auf unerwartete Schwierigkeiten, und die Methode, unter deren Führung wir bisher ohne Hindernisse vorangekommen waren, verweigert jede weitere Hilfe.

\Article{2}

Bevor wir weitergehen, wird es nützlich sein, einige Beispiele unserer Summation zahlenmäßig durchzurechnen; zuvor sollen jedoch einige allgemeine Beobachtungen stehen.

I.\quad Im Fall, dass $n$ eine Primzahl von der Form $4m+1$ ist, seien alle quadratischen Reste von $n$ zwischen $1$ und $\frac12(n-1)$ einschließlich unbestimmt durch $a'$ und alle Nichtreste zwischen denselben Grenzen durch $b'$ dargestellt. Bekanntlich sind dann alle $n-a'$ unter den $a$ und alle $n-b'$ unter den $b$ enthalten. Da alle $a'$, $b'$, $n-a'$, $n-b'$ zusammen offenbar die gesamte Menge der Zahlen $1$, $2$, $3\dotsfour$ $n-1$ ausfüllen, umfassen alle $a'$ zusammen mit allen $n-a'$ sämtliche $a$; ebenso umfassen alle $b'$ zusammen mit allen $n-b'$ sämtliche $b$.\quad Daher ist

\SourceLeaf{23}
\[
\begin{aligned}
\Sigma\cos ak\omega&=\Sigma\cos a'k\omega+\Sigma\cos(n-a')k\omega\\
\Sigma\cos bk\omega&=\Sigma\cos b'k\omega+\Sigma\cos(n-b')k\omega\\
\Sigma\sin ak\omega&=\Sigma\sin a'k\omega+\Sigma\sin(n-a')k\omega\\
\Sigma\sin bk\omega&=\Sigma\sin b'k\omega+\Sigma\sin(n-b')k\omega
\end{aligned}
\]

Da nun $\cos(n-a')k\omega=\cos a'k\omega$, $\cos(n-b')k\omega=\cos b'k\omega$, $\sin(n-a')k\omega=-\sin a'k\omega$, $\sin(n-b')k\omega=-\sin b'k\omega$ gilt, ergibt sich unmittelbar

\[
\begin{aligned}
\Sigma\sin ak\omega&=\Sigma\sin a'k\omega-\Sigma\sin a'k\omega=0\\
\Sigma\sin bk\omega&=\Sigma\sin b'k\omega-\Sigma\sin b'k\omega=0
\end{aligned}
\]

Die Summation der Kosinus nimmt dagegen die Form

\[
\begin{aligned}
\Sigma\cos ak\omega&=2\Sigma\cos a'k\omega\\
\Sigma\cos bk\omega&=2\Sigma\cos b'k\omega
\end{aligned}
\]

an, sodass gelten muss:

\[
\begin{aligned}
1+4\Sigma\cos a'k\omega&=\pm\surd n\\
1+4\Sigma\cos b'k\omega&=\mp\surd n\\
2\Sigma\cos a'k\omega-2\Sigma\cos b'k\omega&=\pm\surd n
\end{aligned}
\]

II.\quad Im Fall, dass $n$ von der Form $4m+3$ ist, ist die Ergänzung jedes Restes $a$ zu $n$ ein Nichtrest und die Ergänzung jedes $b$ ein Rest. Somit stimmen sämtliche $n-a$ mit sämtlichen $b$ und sämtliche $n-b$ mit sämtlichen $a$ überein. Daraus folgt

\[
\begin{aligned}
\Sigma\cos ak\omega=\Sigma\cos(n-b)k\omega=\Sigma\cos bk\omega
\end{aligned}
\]

Da also sämtliche $a$ und $b$ zusammen alle Zahlen $1$, $2$, $3\dotsfour n-1$ ausfüllen und somit $\Sigma\cos ak\omega+\Sigma\cos bk\omega=\cos k\omega+\cos2k\omega+\cos3k\omega+\text{usw.}$ $+\cos(n-1)k\omega=-1$ gilt, ergeben sich die Summationen

\[
\begin{aligned}
\Sigma\cos ak\omega&=-\tfrac12\\
\Sigma\cos bk\omega&=-\tfrac12
\end{aligned}
\]

unmittelbar.\quad Ebenso ist

\[
\begin{aligned}
\Sigma\sin ak\omega=\Sigma\sin(n-b)k\omega=-\Sigma\sin bk\omega
\end{aligned}
\]

\SourceLeaf{24}
woraus ersichtlich ist, wie von den beiden Summationen

\[
\begin{aligned}
2\Sigma\sin ak\omega&=\pm\surd n\\
2\Sigma\sin bk\omega&=\mp\surd n
\end{aligned}
\]

die eine von der anderen abhängt.

\Article{3}

Hier folgen nun die Zahlenrechnungen für einige Beispiele:

I.\quad Für $n=5$ gibt es einen Wert von $a'$, nämlich $a'=1$, und einen Wert von $b'$, nämlich $b'=2$. Es ist

\[
\begin{aligned}
\cos\omega=+0{,}3090169944\qquad\cos2\omega=-0{,}8090169944
\end{aligned}
\]

und somit $1+4\cos\omega=+\surd5$,\quad $1+4\cos2\omega=-\surd5$.

II.\quad Für $n=13$ gibt es drei Werte von $a'$, nämlich $1$, $3$, $4$, und ebenso viele Werte von $b'$, nämlich $2$, $5$, $6$. Daraus berechnen wir

{\small\setlength{\arraycolsep}{2.5pt}\renewcommand{\arraystretch}{1.15}
\[
\begin{array}{rclcrcl}
\cos\omega&=&+0{,}8854560257&\qquad&\cos2\omega&=&+0{,}5680647467\\
\cos3\omega&=&+0{,}1205366803&\qquad&\cos5\omega&=&-0{,}7485107482\\
\cos4\omega&=&-0{,}3546048870&\qquad&\cos6\omega&=&-0{,}9709418174\\\cline{1-3}\cline{5-7}
\text{Summe}&=&+0{,}6513878190&\qquad&\text{Summe}&=&-1{,}1513878189
\end{array}
\]
}

Daher $1+4\Sigma\cos a'\omega=+\surd13$,\quad $1+4\Sigma\cos b'\omega=-\surd13$.

\Needspace{13\baselineskip}III.\quad Für $n=17$ haben wir vier Werte von $a'$, nämlich $1$, $2$, $4$, $8$, und ebenso viele Werte von $b'$, nämlich $3$, $5$, $6$, $7$.\quad Daraus berechnen sich die Kosinus:

{\small\setlength{\arraycolsep}{2.5pt}\renewcommand{\arraystretch}{1.15}
\[
\begin{array}{rclcrcl}
\cos\omega&=&+0{,}9324722294&\qquad&\cos3\omega&=&+0{,}4457383558\\
\cos2\omega&=&+0{,}7390089172&\qquad&\cos5\omega&=&-0{,}2736629901\\
\cos4\omega&=&+0{,}0922683595&\qquad&\cos6\omega&=&-0{,}6026346364\\
\cos8\omega&=&-0{,}9829730997&\qquad&\cos7\omega&=&-0{,}8502171357\\\cline{1-3}\cline{5-7}
\text{Summe}&=&+0{,}7807764064&\qquad&\text{Summe}&=&-1{,}2807764065
\end{array}
\]
}

Daher $1+4\Sigma\cos a'\omega=+\surd17$,\quad $1+4\Sigma\cos b'\omega=-\surd17$.

\SourceLeaf{25}
IV.\quad Für $n=3$ gibt es nur einen Wert von $a$, nämlich $a=1$, dem

\[
\begin{aligned}
\sin\omega=+0{,}8660254038
\end{aligned}
\]

entspricht. Daher $2\sin\omega=+\surd3$.

V.\quad Für $n=7$ gibt es drei Werte von $a$, nämlich $1$, $2$, $4$. Daraus ergeben sich die Sinus:

{\small\setlength{\arraycolsep}{2.5pt}\renewcommand{\arraystretch}{1.15}
\[
\begin{array}{rcll}
\sin\omega&=&+0{,}7818314825&\\
\sin2\omega&=&+0{,}9749279122&\\
\sin4\omega&=&-0{,}4338837391&\\\cline{1-3}
\text{Summe}&=&+1{,}3228756556,&\text{und somit }2\Sigma\sin a\omega=+\surd7.
\end{array}
\]
}

VI.\quad Für $n=11$ sind die Werte von $a$ gleich $1$, $3$, $4$, $5$, $9$. Ihnen entsprechen die Sinus:

{\small\setlength{\arraycolsep}{2.5pt}\renewcommand{\arraystretch}{1.15}
\[
\begin{array}{rcll}
\sin\omega&=&+0{,}5406408175&\\
\sin3\omega&=&+0{,}9898214419&\\
\sin4\omega&=&+0{,}7557495744&\\
\sin5\omega&=&+0{,}2817325568&\\
\sin9\omega&=&-0{,}9096319954&\\\cline{1-3}
\text{Summe}&=&+1{,}6583123952,&\text{und daher }2\Sigma\sin a\omega=+\surd11.
\end{array}
\]
}

\Needspace{19\baselineskip}VII.\quad Für $n=19$ sind die Werte von $a$ gleich $1$, $4$, $5$, $6$, $7$, $9$, $11$, $16$, $17$. Ihnen entsprechen die Sinus:

{\small\setlength{\arraycolsep}{2.5pt}\renewcommand{\arraystretch}{1.15}
\[
\begin{array}{rcll}
\sin\omega&=&+0{,}3246994692&\\
\sin4\omega&=&+0{,}9694002659&\\
\sin5\omega&=&+0{,}9965844930&\\
\sin6\omega&=&+0{,}9157733267&\\
\sin7\omega&=&+0{,}7357239107&\\
\sin9\omega&=&+0{,}1645945903&\\
\sin11\omega&=&-0{,}4759473930&\\
\sin16\omega&=&-0{,}8371664783&\\
\sin17\omega&=&-0{,}6142127127&\\\cline{1-3}
\text{Summe}&=&+2{,}1794494718,&\text{und somit }2\Sigma\sin a\omega=+\surd19.
\end{array}
\]
}

\SourceLeaf{26}
\Article{4}

In allen diesen Beispielen erhält der Wurzelausdruck das positive Vorzeichen. Dasselbe bestätigt sich leicht auch für die größeren Werte $n=23$, $n=29$ usw.; daraus ergibt sich bereits eine hohe Wahrscheinlichkeit, dass dies ebenso allgemein gilt.\quad Ein Beweis dieser Erscheinung lässt sich jedoch nicht aus den an der angeführten Stelle dargelegten Grundsätzen gewinnen; mit vollem Recht ist sie als Gegenstand einer tieferen Untersuchung anzusehen.\quad Diese Abhandlung soll daher einen strengen Beweis dieses höchst eleganten Satzes vorlegen, der zuvor mehrere Jahre lang auf verschiedene Weise vergeblich versucht worden war und schließlich durch ungewöhnliche und recht subtile Überlegungen gelang. Zugleich werden wir den Satz unter Wahrung, oder vielmehr Steigerung, seiner Eleganz zu weit größerer Allgemeinheit führen.\quad Als krönenden Abschluss werden wir schließlich den erstaunlichen und sehr engen Zusammenhang zwischen dieser Summation und einem anderen arithmetischen Satz von höchster Bedeutung zeigen.\quad Wir hoffen, dass diese Untersuchungen den Geometern an sich willkommen sein werden und dass auch die Methoden, durch die wir all dies erreichen konnten und die bei anderen Gelegenheiten ebenfalls nützlich sein können, ihrer Aufmerksamkeit würdig erscheinen werden.

\Article{5}

Unser Beweis geht aus der Betrachtung einer besonderen Art von Reihen hervor, deren Glieder von Ausdrücken wie

\[
\begin{aligned}
\frac{(1-x^m)\,(1-x^{m-1})\,(1-x^{m-2})\dotsfour(1-x^{m-\mu+1})}{(1-x)\,(1-xx)\,(1-x^3)\dotsfive(1-x^\mu)}
\end{aligned}
\]

abhängen. Der Kürze halber bezeichnen wir einen solchen Bruch mit $(m,\mu)$ und schicken zunächst einige allgemeine Beobachtungen über solche Funktionen voraus.

\Needspace{19\baselineskip}I.\quad Ist $m$ eine positive ganze Zahl kleiner als $\mu$, so verschwindet die Funktion $(m,\mu)$ offenbar, da der Zähler den Faktor $1-x^0$ enthält.\quad Für $m=\mu$ stimmen die Faktoren im Zähler in umgekehrter Reihenfolge mit denen im Nenner überein, woraus $(\mu,\mu)=1$ folgt. Für den Fall schließlich, dass $m$ eine positive ganze Zahl größer als $\mu$ ist, gelten die Formeln

\[
\begin{aligned}
(\mu+1,\mu)&=\frac{1-x^{\mu+1}}{1-x}=(\mu+1,1)\\
(\mu+2,\mu)&=\frac{(1-x^{\mu+2})\,(1-x^{\mu+1})}{(1-x)\qquad(1-xx)}=(\mu+2,2)\\
(\mu+3,\mu)&=\frac{(1-x^{\mu+3})(1-x^{\mu+2})(1-x^{\mu+1})}{(1-x)\qquad(1-xx)\qquad(1-x^3)}=(\mu+3,3)\ \text{usw.}
\end{aligned}
\]

\SourceLeaf{27}
oder allgemein

\[
\begin{aligned}
(m,\mu)=(m,m-\mu)
\end{aligned}
\]

II.\quad Ferner bestätigt sich leicht, dass allgemein

\[
\begin{aligned}
(m,\mu+1)=(m-1,\mu+1)+x^{m-\mu-1}(m-1,\mu)
\end{aligned}
\]

gilt. Da ebenso

\[
\begin{aligned}
(m-1,\mu+1)&=(m-2,\mu+1)+x^{m-\mu-2}(m-2,\mu)\\
(m-2,\mu+1)&=(m-3,\mu+1)+x^{m-\mu-3}(m-3,\mu)\\
(m-3,\mu+1)&=(m-4,\mu+1)+x^{m-\mu-4}(m-4,\mu),\ \text{usw.}
\end{aligned}
\]

ist und diese Reihe bis zu

\[
\begin{aligned}
(\mu+2,\mu+1)&=(\mu+1,\mu+1)+x(\mu+1,\mu)\\
&=(\mu,\mu)+x(\mu+1,\mu)
\end{aligned}
\]

fortgesetzt werden kann, sofern $m$ eine positive ganze Zahl größer als $\mu+1$ ist, ergibt sich

\[
\begin{aligned}
(m,\mu+1)&=(\mu,\mu)+x(\mu+1,\mu)+xx(\mu+2,\mu)\\
&\qquad+x^3(\mu+3,\mu)+\text{usw.}\\
&\qquad+x^{m-\mu-1}(m-1,\mu)
\end{aligned}
\]

\Needspace{16\baselineskip}Daraus ist ersichtlich: Ist für einen bestimmten Wert von $\mu$ jede Funktion $(m,\mu)$ bei positivem ganzzahligem $m$ ein Polynom, so muss auch jede Funktion $(m,\mu+1)$ ein Polynom sein.\quad Da diese Voraussetzung für $\mu=1$ zutrifft, gilt sie also auch für $\mu=2$ und damit auch für $\mu=3$ usw. Das heißt, allgemein ist für jeden positiven ganzzahligen Wert von $m$ die Funktion $(m,\mu)$ ein Polynom; gleichbedeutend damit ist das Produkt

\[
\begin{aligned}
(1-x^m)(1-x^{m-1})(1-x^{m-2})\dotsfour(1-x^{m-\mu+1})
\end{aligned}
\]

teilbar durch

\[
\begin{aligned}
(1-x)(1-x^2)(1-x^3)\dotsfour(1-x^\mu)
\end{aligned}
\]

\clearpage\Needspace{14\baselineskip}\Article{6}

Wir betrachten nun zwei Reihen, die beide zu unserem Ziel führen können.\quad Die erste Reihe ist

\SourceLeaf{28}
\[
\begin{aligned}
&1-\frac{1-x^m}{1-x}+\frac{(1-x^m)\,(1-x^{m-1})}{(1-x)\,(1-xx)}\\
&\qquad-\frac{(1-x^m)(1-x^{m-1})(1-x^{m-2})}{(1-x)\quad(1-xx)\quad(1-x^3)}+\text{usw.}
\end{aligned}
\]

oder

\[
\begin{aligned}
1-(m,1)+(m,2)-(m,3)+(m,4)-\text{usw.}
\end{aligned}
\]

die wir der Kürze halber mit $f(x,m)$ bezeichnen.\quad Zunächst ist unmittelbar ersichtlich, dass diese Reihe, sobald $m$ eine positive ganze Zahl ist, nach ihrem Glied an der Stelle $m+1$ (wobei $=\pm1$ gilt) \emph{abbricht} und dass ihre Summe in diesem Fall daher ein endliches Polynom in $x$ sein muss.\quad Ferner ergibt sich aus Art. 5, II, dass allgemein für jeden Wert von $m$

\[
\begin{aligned}
1&=1\\
-(m,1)&=-(m-1,1)-x^{m-1}\\
+(m,2)&=+(m-1,2)+x^{m-2}(m-1,1)\\
-(m,3)&=-(m-1,3)-x^{m-3}(m-1,2),\ \text{usw.}
\end{aligned}
\]

und somit

\[
\begin{aligned}
f(x,m)&=1-x^{m-1}-(1-x^{m-2})(m-1,1)\\
&\qquad+(1-x^{m-3})(m-1,2)\\
&\qquad-(1-x^{m-4})(m-1,3)+\text{usw.}
\end{aligned}
\]

gilt. Offenbar ist aber

\[
\begin{aligned}
(1-x^{m-2})(m-1,1)&=(1-x^{m-1})(m-2,1)\\
(1-x^{m-3})(m-1,2)&=(1-x^{m-1})(m-2,2)\\
(1-x^{m-4})(m-1,3)&=(1-x^{m-1})(m-2,3),\ \text{usw.}
\end{aligned}
\]

woraus wir die folgende Gleichung ableiten:

\[
\begin{aligned}
f(x,m)=(1-x^{m-1})f(x,m-2)\ .\ .\ .\ .\ .\ .\ .\ .\ [1]
\end{aligned}
\]

\clearpage\Needspace{11\baselineskip}\Article{7}

Da für $m=0$ die Gleichung $f(x,m)=1$ gilt, ergibt die soeben gefundene Formel

\[
\begin{aligned}
f(x,2)&=1-x\\
f(x,4)&=(1-x)(1-x^3)\\
f(x,6)&=(1-x)(1-x^3)(1-x^5)\\
f(x,8)&=(1-x)(1-x^3)(1-x^5)(1-x^7),\ \text{usw.}
\end{aligned}
\]

oder allgemein für jeden geraden Wert von $m$:

\[
\begin{aligned}
f(x,m)=(1-x)(1-x^3)(1-x^5)\dotsfour(1-x^{m-1})\ .\ .\ .\ [2].
\end{aligned}
\]

\SourceLeaf{29}
Da hingegen für $m=1$ die Gleichung $f(x,m)=0$ gilt, ist auch

\[
\begin{aligned}
f(x,3)&=0\\
f(x,5)&=0\\
f(x,7)&=0\ \text{usw.}
\end{aligned}
\]

oder allgemein für jeden ungeraden Wert von $m$:

\[
\begin{aligned}
f(x,m)=0
\end{aligned}
\]

Die letztere Summation hätte man im Übrigen schon daraus ableiten können, dass in der Reihe

\[
\begin{aligned}
1-(m,1)+(m,2)-(m,3)+\text{usw.}+(m,m-1)-(m,m)
\end{aligned}
\]

das letzte Glied das erste aufhebt, das vorletzte das zweite usw.

\Article{8}

Für unser Ziel genügt zwar der Fall, dass $m$ eine positive ungerade ganze Zahl ist. Wegen der Eigenart des Gegenstands lohnt es sich aber, auch über die Fälle, in denen $m$ gebrochen oder negativ ist, einiges hinzuzufügen.\quad Unsere Reihe bricht dann offenbar nicht mehr ab, sondern läuft ins Unendliche; außerdem sieht man leicht, dass sie divergent wird, sobald $x$ ein Wert kleiner als $1$ gegeben wird. Ihre Summation muss daher auf Werte von $x$ beschränkt werden, die größer als $1$ sind.

\Needspace{13\baselineskip}Nach Formel [1] in Art. 6 haben wir

\[
\begin{aligned}
f(x,-2)&=\frac{1}{1-\frac1x}\\[8pt]
f(x,-4)&=\frac{1}{1-\frac1x}\pdot\frac{1}{1-\frac1{x^3}}\\[8pt]
f(x,-6)&=\frac{1}{1-\frac1x}\pdot\frac{1}{1-\frac1{x^3}}\pdot\frac{1}{1-\frac1{x^5}},\ \text{usw.}
\end{aligned}
\]

sodass sich der Wert der Funktion $f(x,m)$ auch für einen negativen geraden ganzzahligen Wert von $m$ in endlichen Gliedern angeben lässt.\quad Für die übrigen Werte von $m$ werden wir die Funktion $f(x,m)$ dagegen auf folgende Weise in ein \emph{unendliches Produkt} umwandeln.

Wenn $m$ einem negativen \emph{unendlichen} Wert zustrebt, geht die Funktion $f(x,m)$ über in

\SourceLeaf{30}
\[
\begin{aligned}
&1+\frac1{x-1}+\frac1{x-1}\pdot\frac1{xx-1}\\
&\qquad+\frac1{x-1}\pdot\frac1{xx-1}\pdot\frac1{x^3-1}+\text{usw.}
\end{aligned}
\]

Diese Reihe ist also gleich dem unendlichen Produkt

\[
\begin{aligned}
\frac1{1-\frac1x}\pdot\frac1{1-\frac1{x^3}}\pdot\frac1{1-\frac1{x^5}}\pdot\frac1{1-\frac1{x^7}}\;\text{usw. bis ins Unendliche.}
\end{aligned}
\]

Da ferner allgemein

\[
\begin{aligned}
f(x,m)&=f(x,m-2\lambda)\pdot(1-x^{m-1})(1-x^{m-3})\\
&\qquad(1-x^{m-5})\dotstwo(1-x^{m-2\lambda+1})
\end{aligned}
\]

ist, gilt

\[
\begin{aligned}
f(x,m)&=f(x,-\infty)\pdot(1-x^{m-1})(1-x^{m-3})(1-x^{m-5})\\
&\qquad\text{usw. bis ins Unendliche.}\\
&=\frac{1-x^{m-1}}{1-x^{-1}}\pdot\frac{1-x^{m-3}}{1-x^{-3}}\pdot\frac{1-x^{m-5}}{1-x^{-5}}\pdot\frac{1-x^{m-7}}{1-x^{-7}}\\
&\qquad\text{usw. bis ins Unendliche.}
\end{aligned}
\]

Dessen Faktoren nähern sich im weiteren Verlauf offenbar immer mehr der Einheit.

Besondere Aufmerksamkeit verdient der Fall $m=-1$. Hier ist

\[
\begin{aligned}
f(x,-1)=1+x^{-1}+x^{-3}+x^{-6}+x^{-10}+\text{usw.}
\end{aligned}
\]

\Needspace{9\baselineskip}Diese Reihe ist also gleich dem unendlichen Produkt

\[
\begin{aligned}
\frac{1-x^{-2}}{1-x^{-1}}\pdot\frac{1-x^{-4}}{1-x^{-3}}\pdot\frac{1-x^{-6}}{1-x^{-5}}\;\text{usw.}
\end{aligned}
\]

\Needspace{11\baselineskip}oder, wenn wir $x$ für $x^{-1}$ schreiben,

\[
\begin{aligned}
1+x+x^3+x^6+\text{usw.}&=\frac{1-xx}{1-x}\pdot\frac{1-x^4}{1-x^3}\pdot\frac{1-x^6}{1-x^5}\\
&\qquad\pdot\frac{1-x^8}{1-x^7}\;\text{usw.}
\end{aligned}
\]

Diese Gleichheit zwischen zwei recht verwickelten Ausdrücken, auf die wir bei anderer Gelegenheit zurückkommen werden, ist in der Tat sehr bemerkenswert.
